\section{LLM Prompt Examples}
\label{app:prompts}

This appendix shows representative prompt templates for each evaluation level across all three domains. Each domain uses a domain-specific system prompt and task-specific user prompts. All prompts follow a zero-shot and 3-shot structure.

\subsection{System Prompts}

\begin{small}
\begin{verbatim}
DTI: "You are a pharmaceutical scientist with expertise in drug-target
interactions, assay development, and medicinal chemistry. Provide precise,
evidence-based answers."

CT:  "You are a clinical trial expert with deep knowledge of drug
development, regulatory science, and clinical pharmacology."

PPI: "You are a protein biochemist with expertise in protein-protein
interactions, structural biology, and proteomics experimental methods.
Provide precise, evidence-based answers."
\end{verbatim}
\end{small}

\subsection{L1: Multiple Choice Classification}

\textbf{DTI L1} (4-way: hard negative / conditional negative / methodological negative / dose-time negative):
\begin{small}
\begin{verbatim}
[context_text with evidence description]

Respond with ONLY the letter of the correct answer: A, B, C, or D.
\end{verbatim}
\end{small}

\textbf{CT L1} (5-way: safety / efficacy / enrollment / strategic / other):
\begin{small}
\begin{verbatim}
Based on the clinical trial information below, classify the primary
reason for this trial's failure.

[context_text with trial details]

Categories:
A) Safety -- Trial failed due to drug toxicity, adverse events, or
   safety signals
B) Efficacy -- Trial failed to demonstrate therapeutic benefit vs control
C) Enrollment -- Trial failed to recruit sufficient participants
D) Strategic -- Trial was discontinued for business, strategic, or
   portfolio reasons
E) Other -- Trial failed due to study design flaws, regulatory issues,
   or other reasons

Respond with ONLY a single letter (A, B, C, D, or E).
\end{verbatim}
\end{small}

\textbf{PPI L1} (4-way: direct experimental / systematic screen / computational inference / database score absence):
\begin{small}
\begin{verbatim}
Based on the evidence description below, classify the type and quality
of evidence supporting this protein non-interaction.

[context_text with evidence description]

Categories:
A) Direct experimental -- A specific binding assay (co-IP, pulldown,
   SPR, etc.) found no physical interaction
B) Systematic screen -- A high-throughput binary screen (Y2H, LUMIER,
   etc.) found no interaction
C) Computational inference -- ML analysis of co-fractionation or complex
   data predicts no interaction
D) Database score absence -- Zero or negligible combined interaction
   score across multiple evidence channels

Respond with ONLY a single letter (A, B, C, or D).
\end{verbatim}
\end{small}

\subsection{L2: Structured Extraction}

\textbf{DTI L2} (JSON extraction from abstract):
\begin{small}
\begin{verbatim}
Extract all negative drug-target interaction results from the following
abstract.

Abstract: [abstract_text]

For each negative result found, extract:
- compound: compound/drug name
- target: target protein/gene name
- target_uniprot: UniProt accession (if determinable)
- activity_type: type of measurement (IC50, Ki, Kd, EC50, etc.)
- activity_value: reported value with units
- activity_relation: relation (=, >, <, ~)
- assay_format: biochemical, cell-based, or in vivo
- outcome: inactive, weak, or inconclusive

Also report:
- total_inactive_count: total number of inactive results mentioned
- positive_results_mentioned: true/false

Respond in JSON format.
\end{verbatim}
\end{small}

\textbf{CT L2} (JSON extraction from termination report):
\begin{small}
\begin{verbatim}
Extract structured failure information from the following clinical trial
termination report. Return a JSON object with the fields specified below.

[context_text]

Required JSON fields:
- failure_category: one of [efficacy, safety, pharmacokinetic,
  enrollment, strategic, design, regulatory, other]
- failure_subcategory: specific reason
- affected_system: organ system affected (null if not applicable)
- severity_indicator: one of [mild, moderate, severe, fatal, null]
- quantitative_evidence: true if text mentions specific numbers
- decision_maker: who terminated [sponsor, dsmb, regulatory,
  investigator, null]
- patient_impact: brief description of patient safety impact

Return ONLY valid JSON, no additional text.
\end{verbatim}
\end{small}

\textbf{PPI L2} (JSON extraction of non-interacting pairs):
\begin{small}
\begin{verbatim}
Extract all protein pairs reported as non-interacting from the following
evidence summary. Return a JSON object with the fields specified below.

[context_text]

Required JSON fields:
- non_interacting_pairs: list of objects, each with:
    - protein_1: gene symbol or UniProt accession
    - protein_2: gene symbol or UniProt accession
    - method: experimental method used
    - evidence_strength: one of [strong, moderate, weak]
- total_negative_count: total number of non-interacting pairs mentioned
- positive_interactions_mentioned: true if any positive interactions
  are also mentioned

Return ONLY valid JSON, no additional text.
\end{verbatim}
\end{small}

\subsection{L3: Scientific Reasoning}

\textbf{DTI L3}:
\begin{small}
\begin{verbatim}
[context_text describing compound-target pair and inactivity evidence]

Provide a detailed scientific explanation (3-5 paragraphs) covering:
1. Structural compatibility between compound and target binding site
2. Known selectivity profile and mechanism of action
3. Relevant SAR (structure-activity relationship) data
4. Pharmacological context and therapeutic implications
\end{verbatim}
\end{small}

\textbf{CT L3}:
\begin{small}
\begin{verbatim}
The following clinical trial was confirmed as a FAILURE. Based on the
trial data below, provide a scientific explanation for why this drug
failed in this clinical trial.

[context_text]

Your explanation should address:
1. Mechanism -- What is the drug's mechanism of action and why might
   it be insufficient for this condition?
2. Evidence interpretation -- What do the statistical results tell us?
3. Clinical factors -- Trial design, patient population, or disease
   biology factors?
4. Broader context -- Known challenges of treating this condition?

Provide a thorough explanation in 3-5 paragraphs.
\end{verbatim}
\end{small}

\textbf{PPI L3}:
\begin{small}
\begin{verbatim}
The following two proteins have been experimentally tested and confirmed
to NOT physically interact. Based on the protein information below,
provide a scientific explanation for why they are unlikely to form a
physical interaction.

[context_text]

Your explanation should address:
1. Biological plausibility -- Are there biological reasons (function,
   pathway, localization) that make interaction unlikely?
2. Structural reasoning -- Do domain architectures, binding interfaces,
   or steric factors argue against interaction?
3. Mechanistic completeness -- Are multiple relevant factors considered?
4. Specificity -- Are claims specific to these proteins or generic?

Provide a thorough explanation in 3-5 paragraphs.
\end{verbatim}
\end{small}

\subsection{L4: Tested vs.\ Untested Discrimination}

All three domains use a similar L4 structure:

\begin{small}
\begin{verbatim}
DTI: "Has the following compound-target pair been experimentally
tested for binding activity?"
[compound name + target name]

CT:  "Has the following drug-condition combination ever been tested
in a registered clinical trial?"
[drug name + condition name]

PPI: "Has the following protein pair ever been experimentally tested
for physical interaction?"
[protein A gene symbol + protein B gene symbol]

Response format (all domains):
"On the first line, respond with ONLY 'tested' or 'untested'.
On the second line, provide brief evidence for your answer."
\end{verbatim}
\end{small}

\subsection{Few-Shot Configuration}

For 3-shot prompting, three independent example sets are sampled (seeds 42, 43, 44), and results are reported as mean $\pm$ std across the three sets. Examples are formatted as:

\begin{small}
\begin{verbatim}
Here are examples of [task description]:

[Example 1 context]
Answer: [gold answer]

---

[Example 2 context]
Answer: [gold answer]

---

[Example 3 context]
Answer: [gold answer]

Now [task instruction]:

[Test instance context]
\end{verbatim}
\end{small}
